\subsection{完备滤过模}

命题 5. 设 G 为滤过群，其滤过（G \(_{n}\)）由 G 的不变子群组成。下列条件等价：

\begin{enumerate}
\def\labelenumi{(\alph{enumi})}
\item
  G 是完备拓扑群。
\item
  相伴的 Hausdorff 群 \(\mathbf{G}' = \mathbf{G} / \left(\bigcap_{n\in \mathbf{Z}}\mathbf{G}_n\right)\) 是完备的。
\item
  G 中的每个 Cauchy 序列都收敛。
\end{enumerate}

若 \(\mathbf{G}\) 交换并采用加法记号，则这些条件还等价于下列条件：

\begin{enumerate}
\def\labelenumi{(\alph{enumi})}
\setcounter{enumi}{3}
\tightlist
\item
	元素族 \((x_{\lambda})_{\lambda \in L}\) 是 \(\mathbf{G}'\) 中关于滤子 \(\mathfrak{F}\) 收敛于 0 的元素族，其中该滤子由 \(\mathbf{L}\) 的有限子集的补集构成，并且该元素族在 \(\mathbf{G}'\) 中可求和。
\end{enumerate}

G 上的一个滤子是 Cauchy 滤子（分别为收敛滤子）的充要条件，是它在典范映射 \(G \rightarrow G'\) 下的像是 Cauchy 滤子（分别为收敛滤子）（《一般拓扑学》，第 II 章，§3，第 1 节，命题 4）；首先得到 (a) 与 (b) 的等价性；另一方面，由于 \(G'\) 可度量化，(b) 与 (c) 的等价性由《一般拓扑学》第 IX 章 §2 第 6 节的命题 9 得出。

	现在设 \(G\) 交换。设 \(G'\) 完备，并令 \((x_{\lambda})_{\lambda \in L}\) 为 \(G'\) 中关于 \(\mathfrak{F}\) 收敛于 0 的元素族。对于每个邻域 \(V'\)，它是 \(G'\) 中 0 的邻域且为 \(G'\) 的子群，存在 \(J\)，它是 \(L\) 的有限子集，使得条件 \(\lambda \in L - J\) 蕴含 \(x_{\lambda} \in V'\)；于是有 \(\sum_{\lambda \in H} x_{\lambda} \in V'\)，其中 \(H\) 是 \(L\) 中不与 \(J\) 相交的每个有限子集，这说明元素族 \((x_{\lambda})_{\lambda \in L}\) 可求和（《一般拓扑学》，第 III 章，§ 5，第 2 节，定理 1）。

反之，设条件 (d) 成立，且 \((x_{n})\) 是 \(G'\) 中的 Cauchy 序列；于是族 \((x_{n+1} - x_{n})\) 可求和，特别是通项为 \(x_{n+1} - x_{n}\) 的级数收敛，因而序列 \((x_{n})\) 收敛。

设 G 为滤过群，其滤过 \((\mathbf{G}_{n})\) 由 G 的正规子群组成；商群 \(G/G_{n}\) 是离散的，因而是完备的，因为 \(G_{n}\) 在 G 中是开集。令 \(f_{n}\) 为典范映射 \(G \to G/G_{n}\)，并对 \(m \leqslant n\) 令 \(f_{mn}\) 为典范映射 \(G/G_{n} \to G/G_{m}\)；\((\mathbf{G}/\mathbf{G}_{n}, f_{mn})\) 是以 Z 为指标集的离散群逆系统（《一般拓扑学》，第 III 章，§7，第 3 节）。令 \(\tilde{G}\) 为这个逆系统的逆极限拓扑群，并对所有 n 令 \(g_{n}: \tilde{G} \to G/G_{n}\) 为典范映射；令 \(f: G \to \tilde{G}\) 为映射逆系统 \((f_{n})\) 的逆极限，使得对所有 n 有 \(f_{n} = g_{n} \circ f\)；最后令 j 为 G 到其 Hausdorff 完备化 \(\hat{G}\) 的典范映射；由于 \(G/G_{n}\) 完备，存在唯一的拓扑群同构 i: \(\hat{G} \to \tilde{G}\)，使得 \(f = i \circ j\)（同上，命题 2 的推论 1）；称其为 \(\hat{G}\) 到 \(\tilde{G}\) 的典范同构。

设 H 为另一个滤过群，其滤过 \((\mathbf{H}_{n})\) 由 H 的正规子群组成，并设 \(u: G \to H\) 为与滤过相容的同态（第 4 节）。令 \(\tilde{H} = \lim H/H_{n}\)；对所有 n，u 通过取商定义同态 \(u_{n}: G/G_{n} \xleftarrow{\leftarrow} H/H_{n}\)，且 \(u_{n}\) 显然构成一个映射逆系统；令 \(\tilde{u} = \lim u_{n}\)。此外，令 \(\tilde{H}\) 为 H 的 Hausdorff 完备化，令 \(\hat{u}: \hat{G} \to \hat{H}\) 为 u 经过 Hausdorff 完备化得到的同态（《一般拓扑学》，第 II 章，§ 3，第 7 节，命题 15）。由定义立即可知，若借助典范同构将 \(\tilde{G}\) 识别为 \(\tilde{G}\)、将 \(\tilde{H}\) 识别为 \(\tilde{H}\)，则 \(\hat{u}\) 就识别为 \(\tilde{u}\)。特别地，如果对所有 n，\(u_{n}\) 都是同构，则 \(\hat{u}\) 是拓扑群同构。

\textbf{完备滤过群与环的例}\par

\begin{enumerate}
\def\labelenumi{(\arabic{enumi})}
\item
  设 G 为完备滤过群。G 的每个带诱导滤过的闭子群都是完备的（《一般拓扑学》，第 II 章，§ 3，第 4 节，命题 8）。G 的每个带商滤过的商群都是完备的（《一般拓扑学》，第 IX 章，§ 3，第 1 节，注 1）。
\item
	设 A 为滤过交换环，其滤过记为 \((a_n)_{n\in \mathbf{Z}}\)；令 \(A'\) 为形式幂级数环 \(A[[X_1,\ldots ,X_s]]\)。对于所有 \(e = (e_1,\dots ,e_s)\in \mathbf{N}^s\)，记 \(|e| = \sum_{i\in 1}^{s}e_i\)、\(X^e = \prod_{i\in 1}^{s}X_i^{e_i}\)，于是每个 \(P\in A'\) 都可唯一写成 \(P = \sum_{e\in \mathbf{N}^s}\alpha_{e,P}X^e\)，其中 \(\alpha_{e,P}\in A\)。对于所有 \(n\in \mathbf{Z}\)，令 \(\mathfrak{a}_n'\) 表示由 \(\mathrm{P} \in \mathrm{A}'\) 所成的集合，其中满足 \(\alpha_{e,\mathrm{P}} \in \mathfrak{a}_{n-|e|}\) 对所有 \(e \in \mathbf{N}^s\)；我们证明 \(\mathfrak{a}_n'\) 是 \(\mathrm{A}'\) 的理想。显然，\(\mathfrak{a}_n'\) 是 \(\mathrm{A}'\) 的加法子群；另一方面，若 \(\mathrm{P} \in \mathfrak{a}_n'\) 且 \(\mathrm{Q} \in \mathrm{A}'\)，则对所有 \(e \in \mathbf{N}^s\)，有 \(\alpha_{e,\mathrm{PQ}} = \sum_{e' + e'' = e} \alpha_{e',\mathrm{Q}} \alpha_{e'',\mathrm{P}}\)；由于关系 \(e' + e'' = e\) 蕴含 \(|e''| \leqslant |e|\)，故 \(\mathrm{PQ} \in \mathfrak{a}_n'\)。此外，若 \(\mathrm{Q} \in \mathfrak{a}_m'\)，则当 \(e' + e'' = e\) 时，\(\alpha_{e',\mathrm{Q}} \alpha_{e'',\mathrm{P}} \in \mathfrak{a}_{m-|e'|} \mathfrak{a}_{n-|e''}| \subset \mathfrak{a}_{m+n-|e|}\)，这说明 \((\mathfrak{a}_n')_{n \in \mathbf{Z}}\) 是与 \(\mathrm{A}'\) 的环结构相容的滤过（显然 \(1 \in \mathfrak{a}_0'\)）。今后谈到 \(\mathrm{A}'\) 作为滤过环时，除非另有说明，均指赋予它的滤过 \((\mathfrak{a}_n')\)。显然，\(\bigcap_{n \in \mathbf{Z}} \mathfrak{a}_n'\) 是所有系数均属于 \(\bigcap_{n \in \mathbf{Z}} \mathfrak{a}_n\) 的形式幂级数所成的集合；于是若 A 是 Hausdorff 的，\(\mathrm{A}'\) 也是 Hausdorff 的。若 \(\mathfrak{a}_0 = \mathrm{A}\)，则 \(\mathfrak{a}_0' = \mathrm{A}'\)。
\end{enumerate}

命题 6. 在上述记号下，假设 \(\mathfrak{a}_0 = \mathbf{A}\)，并令 \(h\) 表示映射 \(\mathbf{P} \mapsto (\alpha_{e,\mathbf{P}})_{e \in \mathbb{N}^s}\)。则 \(h\) 是从加法拓扑群 \(\mathbf{A}'\) 到加法拓扑群 \(\mathbf{A}^{\mathbf{N}^s}\) 的同构。多项式环 \(\mathbf{A}[\mathbf{X}_1, \ldots, \mathbf{X}_s]\) 在 \(\mathbf{A}'\) 中稠密；若 \(\mathbf{A}\) 完备，则 \(\mathbf{A}'\) 也完备。

	显然 \(h\) 是双射；\(V_n = h(a_n')\) 是由 \((a_e) \in A^{N^s}\) 所成的集合，其中 \(a_e \in a_{n-|e|}\) 对所有 \(e \in N^s\) 满足 \(|e| \leqslant n\)；由于这样的 \(e\) 只有有限个，\(V_n\) 是 \(A^{N^s}\) 中 0 的邻域。反之，若 \(V\) 是 \(A^{N^s}\) 中 0 的邻域，则存在 \(E\)，它是 \(N^s\) 的有限子集，以及整数 \(v\)，使得条件 \(a_e \in a_v\) 对所有 \(e \in E\) 成立时蕴含 \((a_e) \in V\)；若 \(n\) 是 \(v + |e|\)（\(e \in E\)）所成整数中最大的一个，则 \(h(a_n') \subset V\)，这证明了命题 6 的第一个断言。此外，按上述方式定义 \(n\) 和 \(E\)，\(h(P - \sum_{e \in E} \alpha_{e, P} X^e) \in V\) 对所有 \(P \in A'\) 成立，这说明 \(A[X_1, \ldots, X_s]\) 在 \(A'\) 中稠密。最后一个断言由第一个断言及完备空间的乘积仍完备这一事实得出。

	设 \(m\) 为 \(A\) 的理想，且 \((\mathfrak{a}_n)\) 是 \(m\)-进滤过；若 \(n\) 是 \(A'\) 中由 \(m\) 和 \(X_i\) 生成的理想（\(1 \leqslant i \leqslant s\)），则滤过 \((\mathfrak{a}_n')\) 就是 \(n\)-进滤过。显然，对所有 \(k \geqslant 0\)，\(n^k\) 由元素 \(aX^e\) 生成，其中 \(a \in m^{k-|e|}\)，对所有 \(e \in N^s\) 满足 \(|e| \leqslant k\)，因此 \(n^k \subset \mathfrak{a}_k'\)。反过来证明 \(\mathfrak{a}_k' \subset n^k\)。对所有 \(P \in \mathfrak{a}_k'\)，写成 \(P = P' + P''\)，其中 \(P' = \sum_{|e| < k} \alpha_{e,P} X^e\)，\(P'' = \sum_{|e| \geqslant k} \alpha_{e,P} X^e\)。显然可以写成 \(P'' = \sum_{|e|=k} X^e Q_e\)，其中 \(Q_e\) 是 \(A'\) 中的元素，因此 \(P'' \in n^k\)；另一方面，显然 \(\alpha_{e,P} X^e \in n^k\) 对所有 \(e \in N^s\) 成立，从而 \(P' \in n^k\)。于是 \(n^k = \mathfrak{a}_k'\)。

推论。设 A 为交换环，

\[
\mathbf {A} ^ {\prime} = \mathbf {A} [ [ \mathbf {X} _ {1}, \dots , \mathbf {X} _ {s} ] ]
\]

A 上 s 个不定元的形式幂级数环，n 为 \(\Lambda'\) 的理想，由没有常数项的形式幂级数组成。环 \(\Lambda'\) 在 n-进拓扑下是 Hausdorff 且完备的，而多项式环 \(\mathbf{A}[\mathbf{X}_1,\dots ,\mathbf{X}_s]\) 在 \(\mathbf{A}'\) 中处处稠密。

只需将刚才所述应用于 \(\mathfrak{m} = \{0\}\) 的情形即可。

